\naslov{Enačbe drugega reda}

Zanimajo nas enačbe v dveh spremenljivkah oblike
\[
  Lu = a u_{xx} + 2b u_{xy} + c u_{yy} + d u_x + e u_y + fu = g,
\]
kjer so $a, b, \ldots, g$ dane funkcije.
Vpeljimo kvadratno formo, ki je prirejena glavnemu delu operatorja $L$:
\[
  \sigma_{(x,y)}(\xi, \eta)
  = a(x,y) \xi^2 + 2 b(x,y) \xi \eta + c(x,y) \eta^2.
\]
Ob izbrani točki $(x,y) \in \R^2$ je to kvadratna forma na $\R^2$.
Imenuje se \pojem{simbol} operatorja $L$.
Sedaj izračunamo
\[
  \sigma_{(x,y)}(\xi, \eta)
  =
  \sk{
	\begin{bmatrix}
	  a(x,y) & b(x,y) \\
	  b(x,y) & c(x,y)
	\end{bmatrix}
	\cdot
	\begin{bmatrix}
	  \xi \\ \eta
	\end{bmatrix},
	\begin{bmatrix}
	  \xi \\ \eta
	\end{bmatrix}
  }_{\R^2},
\]
in označimo zgornjo matriko z $A(x,y)$.
Poleg tega označimo $\delta = b^2 - ac$.

\begin{definicija}
  Operator $L$ je v točki $(x,y)$
  \begin{itemize}
  \item \pojem{eliptičen}, če je $\delta < 0$,
  \item \pojem{hiperboličen}, če je $\delta > 0$,
  \item \pojem{paraboličen}, če je $\delta = 0$.
  \end{itemize}
\end{definicija}

\vprasanje{Kakšne vrste diferencialnih operatorjev drugega reda poznamo?}

Imejmo $L_0 u = a u_{xx} + 2b u_{xy} + c u_{yy}$ in funkciji $\xi, \eta : \R^2
\to \R$.
Naj bo $w(\xi, \eta) = w(\xi(x,y), \eta(x,y)) = u(x,y)$ pretvorba v nove
koordinate.
Z verižnim pravilom pridemo do
\[
  Lu = \tilde{L} w
  = A w_{\xi \xi} + 2B w_{\xi \eta} + C w_{\eta \eta} + D w_{\xi} + E w_{\eta} + Fw
\]
za
\begin{gather*}
  A = a \xi_x^2 + 2b \xi_x \xi_y + c \xi_y^2 \\
  B = a \xi_x \eta_x + b (\xi_x \eta_y + \xi_y \eta_x) + c \xi_y \eta_y \\
  C = a \eta_x^2 + 2b \eta_x \eta_y + c \eta_y^2.
\end{gather*}